\par\Needspace{4\baselineskip}
\begin{center}
\hypertarget{native-062-article045}{}\pdfbookmark[2]{45.}{article-045}45.
\end{center}
Der zweite Teil der ersten diesem Werk beigefügten Tafel betrifft, wie schon oben erwähnt, die hyperbolische Bewegung. Für das Argument $A$, das beiden Teilen der Tafel gemeinsam ist, gibt er den Logarithmus von $B$ und die Größe $C$ auf sieben Dezimalstellen an, wobei die führenden Ziffern weggelassen sind, und die Größe $T$ zunächst auf fünf, dann auf sechs Stellen. Dieser Teil reicht ebenso wie der erste bis $A=0{,}300$. Dem entsprechen $T=0{,}241207$, $u=2{,}930$ oder $=0{,}341$ und $F=\pm52^0\,19'$. Eine weitere Ausdehnung wäre überflüssig gewesen (Art. 36).

Wir geben nun den Rechengang sowohl für die Bestimmung der Zeit aus der wahren Anomalie als auch für die Bestimmung der wahren Anomalie aus der Zeit an. In der ersten Aufgabe erhält man $T$ nach der Formel $T=\frac{e-1}{e+1}\operatorname{tang}\tfrac{1}{2}v^2$. Zu $T$ liefert unsere Tafel $\log B$ und $C$, woraus $A=\frac{(1+C)T}{1-\tfrac{4}{5}T}$ folgt. Daraus findet man schließlich nach Formel 2) des vorigen Artikels $t$. In der zweiten Aufgabe berechnet man zunächst die Logarithmen der Konstanten
\[
\begin{aligned}
\alpha&=\frac{75k\BareRoot{(\tfrac{1}{5}+\tfrac{9}{5}e)}}{2q^{\frac{3}{2}}}\mbox{\HistoricalDisplayNoteMark}\\[.7em]
\beta&=\frac{5e-5}{1+9e}\\[.7em]
\gamma&=\BareRoot{\frac{5e+5}{1+9e}}\mathord{\raisebox{.35ex}{.}}
\end{aligned}
\]
\noindent\HistoricalDisplayNoteText{Handschriftliche Bemerkung zu einer Tafel, die \textsc{Gauss} 1816 für seinen eigenen Gebrauch berechnete. Die Tafel ist noch ungedruckt und wird in diesem Band erneut besprochen:] in unserer Tafel $\alpha=\frac{k\BareRoot{(\tfrac{1}{5}+\tfrac{9}{5}e)}}{q^{\frac{3}{2}}}\mathord{\raisebox{.35ex}{.}}$}
\SourceLeaf{063}Dann bestimmt man $A$ aus $t$ genau wie bei der elliptischen Bewegung. Der mittleren Bewegung $\frac{\alpha t}{B}$ entspricht in der \textsc{Barker}schen Tafel die wahre Anomalie $w$, und es gilt $A=\beta\operatorname{tang}\tfrac{1}{2}w^2$. Zunächst ermittelt man einen Näherungswert für $A$, wobei man den Faktor $B$ vernachlässigt oder, falls die nötigen Hilfsmittel vorliegen, schätzt. Unsere Tafel liefert dann einen Näherungswert für $B$, mit dem man die Rechnung wiederholt. Der so erhaltene neue Wert von $B$ wird kaum je einer merklichen Korrektur bedürfen, sodass man die Rechnung kein weiteres Mal wiederholen muss. Nach der Korrektur von $A$ entnimmt man der Tafel $C$. Damit ergibt sich
\[
\begin{aligned}
\operatorname{tang}\tfrac{1}{2}v&=\frac{\gamma\operatorname{tang}\tfrac{1}{2}w}{\BareRoot{(1+\tfrac{4}{5}A+C)}},\\[.7em]
r&=\frac{(1+\tfrac{4}{5}A+C)q}{(1-\tfrac{1}{5}A+C)\cos\tfrac{1}{2}v^2}\mathord{\raisebox{.35ex}{.}}
\end{aligned}
\]
Daraus geht hervor, dass die Formeln für die elliptische und die hyperbolische Bewegung genau übereinstimmen, wenn wir $\beta$, $A$ und $T$ bei der hyperbolischen Bewegung als negative Größen behandeln.
